% math0 manuscript. Build with: latexmk -pdf paper.tex
\documentclass[11pt]{amsart}
\usepackage[margin=1.1in]{geometry}
\usepackage{amsmath,amssymb,amsthm,mathtools}
\usepackage{url}
\usepackage{hyperref}
\usepackage{xcolor}
\usepackage{listings}
\usepackage{newunicodechar}

\lstdefinelanguage{lean}{
  keywords={theorem,lemma,def,abbrev,example,fun,by,have,show,from,let,in,if,then,else,match,with,
    namespace,end,open,import,section,variable,instance,structure,where,calc,obtain,intro,exact,
    refine,rcases,rintro,simp,omega,decide,norm_num,linarith,nlinarith,positivity,induction,cases,
    private,noncomputable},
  sensitive=true, comment=[l]{--}, morecomment=[s]{/-}{-/}, string=[b]"}
\lstset{language=lean, basicstyle=\ttfamily\footnotesize, columns=fullflexible, keepspaces=true, breaklines=true,
  keywordstyle=\color{blue!60!black}\bfseries, commentstyle=\color{gray}, frame=single,
  extendedchars=true, inputencoding=utf8,
  literate={ℕ}{{$\mathbb{N}$}}1 {∀}{{$\forall$}}1 {→}{{$\to$}}1 {≤}{{$\le$}}1 {∈}{{$\in$}}1
    {∧}{{$\wedge$}}1 {∑ᶠ}{{$\sum^{\mathrm{f}}$}}2 {·}{{$\cdot$}}1}
% Unicode commonly found in Lean sources.
\newunicodechar{ℕ}{\ensuremath{\mathbb{N}}}
\newunicodechar{ℤ}{\ensuremath{\mathbb{Z}}}
\newunicodechar{ℚ}{\ensuremath{\mathbb{Q}}}
\newunicodechar{ℝ}{\ensuremath{\mathbb{R}}}
\newunicodechar{ℂ}{\ensuremath{\mathbb{C}}}
\newunicodechar{∀}{\ensuremath{\forall}}
\newunicodechar{∃}{\ensuremath{\exists}}
\newunicodechar{→}{\ensuremath{\to}}
\newunicodechar{←}{\ensuremath{\leftarrow}}
\newunicodechar{↔}{\ensuremath{\leftrightarrow}}
\newunicodechar{↦}{\ensuremath{\mapsto}}
\newunicodechar{¬}{\ensuremath{\neg}}
\newunicodechar{∧}{\ensuremath{\wedge}}
\newunicodechar{∨}{\ensuremath{\vee}}
\newunicodechar{≤}{\ensuremath{\le}}
\newunicodechar{≥}{\ensuremath{\ge}}
\newunicodechar{≠}{\ensuremath{\ne}}
\newunicodechar{∈}{\ensuremath{\in}}
\newunicodechar{∉}{\ensuremath{\notin}}
\newunicodechar{⊆}{\ensuremath{\subseteq}}
\newunicodechar{∩}{\ensuremath{\cap}}
\newunicodechar{∪}{\ensuremath{\cup}}
\newunicodechar{∑}{\ensuremath{\sum}}
\newunicodechar{∏}{\ensuremath{\prod}}
\newunicodechar{⟨}{\ensuremath{\langle}}
\newunicodechar{⟩}{\ensuremath{\rangle}}
\newunicodechar{·}{\ensuremath{\cdot}}
\newunicodechar{∘}{\ensuremath{\circ}}
\newunicodechar{⁻}{\ensuremath{^{-}}}
\newunicodechar{¹}{\ensuremath{^{1}}}
\newunicodechar{²}{\ensuremath{^{2}}}
\newunicodechar{₀}{\ensuremath{_{0}}}
\newunicodechar{₁}{\ensuremath{_{1}}}
\newunicodechar{₂}{\ensuremath{_{2}}}
\newunicodechar{α}{\ensuremath{\alpha}}
\newunicodechar{β}{\ensuremath{\beta}}
\newunicodechar{ε}{\ensuremath{\varepsilon}}
\newunicodechar{δ}{\ensuremath{\delta}}
\newunicodechar{π}{\ensuremath{\pi}}
\newunicodechar{σ}{\ensuremath{\sigma}}
\newunicodechar{φ}{\ensuremath{\varphi}}
\newunicodechar{∞}{\ensuremath{\infty}}
\newunicodechar{⌊}{\ensuremath{\lfloor}}
\newunicodechar{⌋}{\ensuremath{\rfloor}}
\newunicodechar{⌈}{\ensuremath{\lceil}}
\newunicodechar{⌉}{\ensuremath{\rceil}}
\newunicodechar{∣}{\ensuremath{\mid}}
\newunicodechar{≡}{\ensuremath{\equiv}}
\newunicodechar{×}{\ensuremath{\times}}
\newunicodechar{⧸}{\ensuremath{/}}
\newunicodechar{‖}{\ensuremath{\|}}
\newunicodechar{√}{\ensuremath{\surd}}
\newunicodechar{∅}{\ensuremath{\emptyset}}
\newunicodechar{⦃}{\ensuremath{\{\!|}}
\newunicodechar{⦄}{\ensuremath{|\!\}}}
\newunicodechar{𝓝}{\ensuremath{\mathcal{N}}}
\newunicodechar{≫}{\ensuremath{\gg}}
\newunicodechar{≪}{\ensuremath{\ll}}
\newunicodechar{ᶠ}{\ensuremath{^{\mathrm{f}}}}

\theoremstyle{plain}
\newtheorem{theorem}{Theorem}
\newtheorem{lemma}[theorem]{Lemma}
\newtheorem{proposition}[theorem]{Proposition}
\newtheorem{corollary}[theorem]{Corollary}
\theoremstyle{definition}
\newtheorem{definition}[theorem]{Definition}
\newtheorem{remark}[theorem]{Remark}

\emergencystretch=3em
\allowdisplaybreaks
% Lean identifiers in running text (breakable at dots and underscores).
\expandafter\def\expandafter\UrlBreaks\expandafter{\UrlBreaks\do\_\do\.}
\DeclareUrlCommand\decl{\def\UrlFont{\normalfont\ttfamily}}
\newcommand{\G}{\mathcal{G}}
\newcommand{\lb}{\char123}
\newcommand{\rb}{\char125}

\title[Column sums of OEIS A388724]{Strictly superdiagonal reversed partitions with largest part $k$ are counted by $\mathrm{A000245}(k-1)$: a formal proof of Wiseman's conjecture for OEIS A388724}
\author{math0}
\address{Produced by Claude (an AI model developed by Anthropic) in the math0 project}
\date{October 25, 2026}

\begin{document}
\begin{abstract}
Let $T(n,k)$ be the number of integer partitions of $n$ whose parts, listed in weakly increasing order as $y_1\le\dots\le y_m$, satisfy $y_i>i$ for all $i$ and have largest part $k$ (OEIS A388724).
We prove that $\sum_{n}T(n,k)=\mathrm{A000245}(k-1)$ for every $k\ge1$, where $\mathrm{A000245}(0)=0$ and $\mathrm{A000245}(K)=3(2K)!/((K+2)!\,(K-1)!)$ for $K\ge1$; this is the conjecture ``Column sums are A000245'' posted by Gus Wiseman in the entry in September 2025.
The proof counts such sequences with entries in a prescribed interval by splitting off the first entry, which produces Catalan's triangle, and evaluates one entry of that triangle by the ballot formula.
The conjecture is still listed as open in the OEIS, but it follows quickly from bijections of Archibald, Blecher, Elizalde and Knopfmacher (2025) and classical Dyck-path counts, so we classify our result as the first formal proof of a known result.
It is verified in Lean~4 against the frozen, fidelity-reviewed statement that math0 formalized from the OEIS entry, whose definitions are checked by the kernel against all $78$ listed terms.
\end{abstract}
\maketitle

\section{Introduction}

A finite sequence $(y_1,\dots,y_m)$ is \emph{strictly superdiagonal} if $y_i>i$ for $i=1,\dots,m$. A \emph{reversed partition} of $n$ is a partition of $n$ with its parts listed in weakly increasing order. The OEIS entry A388724 \cite{A388724}, submitted by Gus Wiseman on September~24, 2025, is the triangle
\[
  T(n,k)=\#\{\text{strictly superdiagonal reversed partitions of $n$ with maximum $k$}\},\qquad 0\le k\le n .
\]
The entry contains the line ``Conjecture: Column sums are A000245.'', where A000245 \cite{A000245} is $0,1,3,9,28,90,297,\dots$, given by $3(2n)!/((n+2)!(n-1)!)$ for $n\ge1$. Since a part $y_m=k$ forces $m<k$ and hence $n\le k(k-1)$, every column has finite support. The column sums for $k=0,1,2,\dots$ are $0,0,1,3,9,28,\dots$, so the conjecture says that column $k$ sums to $\mathrm{A000245}(k-1)$ for $k\ge1$. Equivalently: \emph{the number of weakly increasing sequences of integers $(y_1,\dots,y_m)$, $m\ge1$, with $y_i>i$ for all $i$ and $y_m=k$, is $\mathrm{A000245}(k-1)$.} We prove this for all $k\ge1$ (Theorem~\ref{thm:main}), with a Lean proof of the full statement.

\subsection*{Prior work and status}
\begin{itemize}
\item \emph{The OEIS entry.} In revision \#12 of A388724 (October~6, 2025), which is still the latest revision in October 2026, the conjecture is stated and no proof is recorded. The entry also contains a generating function for the triangle, a PARI program and a b-file of rows $0$ to $100$ by John Tyler Rascoe (October~4, 2025), but no proof of the conjecture.
\item \emph{A short derivation from published results.} Archibald, Blecher, Elizalde and Knopfmacher \cite{ABEK} study \emph{superdiagonal} weakly increasing partitions, $y_i\ge i$. Subtracting $1$ from every part maps the strictly superdiagonal reversed partitions with largest part $k$ bijectively onto the superdiagonal weakly increasing partitions with largest part $k-1$ (all parts stay positive since $y_i-1\ge i$). In the proof of their Theorem~2.3, a superdiagonal weakly increasing partition with largest part $m$ corresponds to a pair $(P,h)$, where $P$ is a Dyck path of semilength $m$ and $0\le h<\ell(P)$, with $\ell(P)$ the length of the final run of east steps of $P$. So the column sum is $\sum_{P}\ell(P)$ over Dyck paths of semilength $k-1$, and this total is classically $\mathrm{A000245}(k-1)$: reversing a path turns $\ell(P)$ into the height of its first peak, which has the same distribution as the number of returns to the axis (a comment of E.~Deutsch in A033184 \cite{A033184}; see also his paper \cite{Deutsch}), and by a comment of C.~Homberger in A000245 \cite{A000245} the total number of returns over all Dyck paths of semilength $n$ is $\mathrm{A000245}(n)$. Alternatively, in the proof of their Theorem~2.2, the partitions with $r$ parts and largest part $s$ correspond to ballot paths ending at $(r-1,s)$; these are counted by entries of Catalan's triangle A009766 \cite{A009766}, and summing over $r$ gives $C_{s+1}-C_s=\mathrm{A000245}(s)$, where $C_s$ is the $s$-th Catalan number. This derivation was pointed out by the math0 literature referee; we have not formalized it. The paper \cite{ABEK} appeared on April~7, 2025, before the conjecture was posted.
\item The math0 referees found no explicit statement or proof of the column-sum identity in the OEIS, arXiv, MathOverflow, Mathematics Stack Exchange or elsewhere on the web.
\end{itemize}
As the identity is in substance a corollary of known results, we classify our result as the \emph{first formal proof of a known result}. Our proof is self-contained: it uses only the definition of $T$, a recursion on the first entry, and binomial identities.

\section{Statement}\label{sec:statement}

The entry A388724 \cite{A388724} (offset $0$, keywords \texttt{nonn,tabl}) has the name ``Triangle read by rows where T(n,k) is the number of strictly superdiagonal reversed integer partitions of n with maximum k.'', the comment ``A sequence (y\_1, ..., y\_k) is strictly superdiagonal iff y\_i \textgreater{} i for all i = 1..k.'', and, in its FORMULA section, verbatim:
\begin{quote}\raggedright
\texttt{Conjecture: Column sums are A000245.}
\end{quote}
The line carries no signature of its own; the revision history of the entry shows that it was added by the author of the entry, Gus Wiseman, in revision \#2 (September~24, 2025). The entry A000245 \cite{A000245} (N.~J.~A. Sloane, offset $0$) has the name ``a(n) = 3*(2*n)!/((n+2)!*(n-1)!).'' and data $0,1,3,9,28,90,\dots$.

\begin{theorem}\label{thm:main}
For $n,k\ge0$ let $T(n,k)$ be the number of finite sequences $y=(y_1,\dots,y_m)$, $m\ge0$, of positive integers with $y_1\le\dots\le y_m$, $y_1+\dots+y_m=n$, $y_i>i$ for $i=1,\dots,m$, and $\max_i y_i=k$ (so $m\ge1$). Let $\mathrm{A000245}(0)=0$ and $\mathrm{A000245}(K)=\dfrac{3\,(2K)!}{(K+2)!\,(K-1)!}$ for $K\ge1$. Then for every $k\ge1$,
\[
  \sum_{n\ge0}T(n,k)=\mathrm{A000245}(k-1),
\]
where the sum is finite since $T(n,k)=0$ for $n>k(k-1)$.
\end{theorem}

\subsection*{Lean formalization}
The statement was formalized by math0 in \path{lean/Math0/Statements/OEIS/A388724.lean}, reviewed for fidelity, and frozen (SHA-256 \texttt{62a9cf37\ldots 798e04}) before any proof attempt. Its definitions and the statement read as follows (verbatim; we omit the import, comments, blank lines, the line \texttt{open Finset}, a \texttt{DecidablePred} instance for \decl{StrictlySuperdiagonal}, all theorems, and the definitions \decl{rowOf}, \decl{a}, \decl{gen} and \decl{Tc} used for the data checks described below):
\begin{lstlisting}
namespace Math0.Statements.A388724
def IsRevPartition (n : ℕ) (y : List ℕ) : Prop :=
  (∀ x ∈ y, 0 < x) ∧ y.Pairwise (· ≤ ·) ∧ y.sum = n
def StrictlySuperdiagonal (y : List ℕ) : Prop := ∀ i (h : i < y.length), i + 1 < y[i]
noncomputable def T (n k : ℕ) : ℕ :=
  {y : List ℕ | IsRevPartition n y ∧ StrictlySuperdiagonal y ∧ y.max? = some k}.ncard
def A000245 (n : ℕ) : ℕ :=
  if n = 0 then 0 else 3 * (2 * n).factorial / ((n + 2).factorial * (n - 1).factorial)
def statement : Prop := ∀ k : ℕ, 1 ≤ k → ∑ᶠ n, T n k = A000245 (k - 1)
end Math0.Statements.A388724
\end{lstlisting}

\emph{Encoding.} A partition of $n$ is represented by the unique list of its parts in weakly increasing order (\texttt{List.Pairwise (· ≤ ·)}), so \decl{IsRevPartition} is the set of reversed partitions of $n$; the empty list is the partition of $0$. Lists are $0$-indexed, so \texttt{i + 1 < y[i]} for $i<m$ is $y_i>i$ for $1\le i\le m$. ``Maximum $k$'' is \texttt{y.max?\ = some k}; the empty list has no maximum and is never counted, in agreement with the listed $T(0,0)=0$ and with the Mathematica program of the entry. \texttt{Set.ncard} is the cardinality of a finite set ($0$ for an infinite one); this set is finite by \decl{T_eq_Tc} below. In \decl{A000245} the division is exact, and the value at $0$ is the listed term $0$. The column sum is the \texttt{finsum} $\sum^{\mathrm f}_n T(n,k)$, which is the ordinary sum when the support is finite (and $0$ otherwise); the statement file proves \decl{T_eq_zero_of_lt}: $T(n,k)=0$ for $n>k(k-1)$, and \decl{colSum_eq}: $\sum^{\mathrm f}_nT(n,k)=\sum_{n=0}^{k(k-1)}T(n,k)$, so the sum is an honest finite sum. Column $0$ is identically zero, and the alignment $k\mapsto k-1$ is the only one compatible with the listed values ($\mathrm{A000245}(1)=1\neq0=\sum_nT(n,1)$); the statement claims nothing about column $0$.

\emph{Data checks.} The statement file defines a fuel-bounded generator \decl{gen} of weakly increasing lists with given sum and lower bound, proves it correct and complete (\decl{mem_gen}) and free of repetitions (\decl{nodup_gen}), and deduces \decl{T_eq_Tc}: $T(n,k)$ equals the length \decl{Tc} of a filtered list. Through this lemma, \texttt{decide +kernel} checks that the triangle read by rows (\decl{a}) reproduces all $78$ listed terms of A388724 (rows $0$ to $11$; \decl{a_data}), that \decl{A000245} reproduces the first $28$ listed terms of A000245 (\decl{A000245_data}), and that the statement holds for $k=1,\dots,5$ (\decl{statement_check}), the last as a transcription check only.

\section{Proof}\label{sec:proof}

The proof follows \path{lean/Math0/Frontier/OEIS/A388724.lean}; the names of the corresponding Lean declarations are given in parentheses. All sequences below are finite sequences of nonnegative integers, the empty sequence included, and we write $x\in y$ if $x$ is an entry of $y$.

\begin{definition}
For $k,\ell\ge0$ let $\G(k,\ell)$ be the set of weakly increasing sequences $y=(y_1,\dots,y_m)$, $m\ge0$, with $y_i>i$ for all $i$ and $\ell\le y_i\le k$ for all $i$ (\decl{Good}~$\ell$~$k$), and let $N(k,\ell)=|\G(k,\ell)|$ (\decl{N}).
\end{definition}

Since $y_m>m$, every nonempty $y\in\G(k,\ell)$ has length $m<k$, and its entries are at most $k$, so $\G(k,\ell)$ is finite. In Lean, strict superdiagonality is expressed through the recursive predicate $\operatorname{SD}_c$ (\decl{SD}): $\operatorname{SD}_c()$ holds and $\operatorname{SD}_c(x,t)$ means $x>c$ and $\operatorname{SD}_{c+1}(t)$; thus $\operatorname{SD}_c(y)$ says that $y_i>i-1+c$ for all $i$ (\decl{SD_iff}), and $\operatorname{SD}_1$ is strict superdiagonality.

\subsection{Reduction to counting sequences}

\begin{lemma}[\decl{colSum_card}]\label{lem:colsum}
For every $k\ge0$, $\sum_nT(n,k)=\bigl|\{y\in\G(k,0): k\in y\}\bigr|$.
\end{lemma}

\begin{proof}
Let $F_k=\{y\in\G(k,0):k\in y\}$. For a sequence $y$ and $n\ge0$, the conditions defining $T(n,k)$ hold if and only if $y\in F_k$ and $y_1+\dots+y_m=n$: positivity of the parts follows from $y_i>i\ge1$, and $\max y=k$ means that $k\in y$ and all entries are at most $k$. So $T(n,k)$ is the size of the fibre over $n$ of the map $y\mapsto\sum_iy_i$ on the finite set $F_k$. The support of $n\mapsto T(n,k)$ is contained in the image of this map, and the sum of the sizes of all fibres is $|F_k|$.
\end{proof}

\subsection{Counting by the first entry}
For a sequence $t=(t_1,\dots,t_r)$ write $t+1=(t_1+1,\dots,t_r+1)$, and write $j\cdot s$ for the sequence with first entry $j$ followed by $s$.

\begin{lemma}[\decl{enum}, \decl{mem_enum}, \decl{nodup_enum}, \decl{N_succ}]\label{lem:enum}
$\G(0,\ell)=\{()\}$ for every $\ell$. For $k\ge0$, the map $(j,t)\mapsto j\cdot(t+1)$ is a bijection from
\[
  \{(j,t):\ \max(\ell,2)\le j\le k+1,\ t\in\G(k,j-1)\}
\]
onto the set of nonempty sequences in $\G(k+1,\ell)$. Consequently $N(0,\ell)=1$ and
\begin{equation}\label{eq:Nrec}
  N(k+1,\ell)=1+\sum_{j=\max(\ell,2)}^{k+1}N(k,j-1).
\end{equation}
\end{lemma}

\begin{proof}
A nonempty $y\in\G(0,\ell)$ would satisfy $1<y_1\le0$. Let $y=j\cdot s$ be nonempty, $s=(s_1,\dots,s_r)$, and suppose all $s_i\ge1$, so that $s=t+1$ with $t_i=s_i-1$. Then:
\begin{itemize}
\item $y$ is strictly superdiagonal if and only if $j>1$ and $s_i>i+1$ for all $i$, i.e.\ $t_i>i$; in Lean, $\operatorname{SD}_{c+1}(t+1)\iff\operatorname{SD}_c(t)$ (\decl{SD_map});
\item $y$ is weakly increasing if and only if $t$ is weakly increasing and $t_i\ge j-1$ for all $i$;
\item given these, $\ell\le y_i\le k+1$ for all $i$ if and only if $\ell\le j\le k+1$ and $t_i\le k$ for all $i$ (the bound $t_i+1\ge j\ge\ell$ is automatic).
\end{itemize}
Hence $j\cdot(t+1)\in\G(k+1,\ell)$ if and only if $\max(\ell,2)\le j\le k+1$ and $t\in\G(k,j-1)$. Conversely every nonempty $y\in\G(k+1,\ell)$ has the form $j\cdot s$ with $s_i\ge j\ge2$, so it is $j\cdot(t+1)$ with $t=s-1$ (\decl{map_sub_add}). The map is injective, since $j$ is the first entry and $t\mapsto t+1$ is injective.

In Lean, \decl{enum}~$k$~$\ell$ is \texttt{[[]]} for $k=0$, and for $k+1$ it is the empty list followed, for $j=\max(\ell,2),\dots,k+1$, by the lists $j\cdot(t+1)$, $t\in{}$\decl{enum}~$k$~$(j-1)$. By the argument above its members are the elements of $\G(k,\ell)$ (\decl{mem_enum}) and it has no repetitions (\decl{nodup_enum}); $N(k,\ell)$ is its length, and \eqref{eq:Nrec} is read off from its shape (\decl{N_succ}).
\end{proof}

\begin{corollary}\label{cor:N}
For all $k,\ell\ge0$:
\begin{enumerate}
\item[(a)] $N(k+1,\ell)=N(k+1,2)$ if $\ell\le2$, and $N(k,1)=N(k,0)$ \quad(\decl{N_low}, \decl{N_one});
\item[(b)] $N(k+1,\ell)=1$ if $\ell\ge k+2$ \quad(\decl{N_top});
\item[(c)] $N(k+1,\ell)=N(k+1,\ell+1)+N(k,\ell-1)$ if $2\le\ell\le k+1$ \quad(\decl{N_rec}).
\end{enumerate}
\end{corollary}

\begin{proof}
(a) If $\ell\le2$ then $\max(\ell,2)=2$ in \eqref{eq:Nrec}; for $k=0$ in the second claim, $N(0,1)=N(0,0)=1$. (b) The sum in \eqref{eq:Nrec} is empty. (c) For $\ell\ge2$ the sum in \eqref{eq:Nrec} runs over $\ell\le j\le k+1$; splitting off the term $j=\ell$ leaves the sum for $\ell+1$.
\end{proof}

\subsection{Catalan's triangle}
For $k,s\ge0$ put $A(k,s)=N(k,k+1-s)$ (\decl{A}; truncated subtraction, but we only use $s\le k$). Thus $A(k,s)$ counts the sequences in $\G(k,k+1-s)$, whose entries lie in $[k+1-s,k]$.

\begin{lemma}\label{lem:A}
For all $k,s\ge0$:
\begin{enumerate}
\item[(a)] $A(k,0)=1$ \quad(\decl{A_zero});
\item[(b)] $A(k+1,s+1)=A(k+1,s)+A(k,s+1)$ if $s+1\le k$ \quad(\decl{A_rec});
\item[(c)] $A(k+1,k+1)=A(k+1,k)$ \quad(\decl{A_diag}).
\end{enumerate}
\end{lemma}

\begin{proof}
(a) $A(0,0)=N(0,1)=1$, and $A(k+1,0)=N(k+1,k+2)=1$ by Corollary~\ref{cor:N}(b). (b) Apply Corollary~\ref{cor:N}(c) with $\ell=k+1-s$, which satisfies $2\le\ell\le k+1$: it reads $N(k+1,k+1-s)=N(k+1,k+2-s)+N(k,k-s)$, which is (b). (c) $A(k+1,k+1)=N(k+1,1)=N(k+1,2)=A(k+1,k)$ by Corollary~\ref{cor:N}(a).
\end{proof}

\begin{proposition}[\decl{A_closed}]\label{prop:ballot}
For all $n,m\ge0$ with $m+1\le n$,
\[
  A(n,m+1)+\binom{n+m+1}{m}=\binom{n+m+1}{m+1}.
\]
\end{proposition}

So $A(n,s)=\binom{n+s}{s}-\binom{n+s}{s-1}$ for $0\le s\le n$ (with $s=0$ from Lemma~\ref{lem:A}(a)), which are the entries of Catalan's triangle A009766 \cite{A009766}.

\begin{proof}
Induction on $n$; for $n=0$ there is nothing to prove. Assume the claim for $n$ (for all admissible $m$) and prove it for $n+1$ by induction on $m$, where $m+1\le n+1$.

\emph{Case $m=0$}: we need $A(n+1,1)=n+1$. If $n=0$, $A(1,1)=A(1,0)=1$ by Lemma~\ref{lem:A}(c),(a). If $n\ge1$, Lemma~\ref{lem:A}(b) with $s=0$ gives $A(n+1,1)=A(n+1,0)+A(n,1)=1+n$, using the outer hypothesis $A(n,1)+1=n+1$.

\emph{Case $m+1$ with $m+2\le n$} (generic step). By Lemma~\ref{lem:A}(b), $A(n+1,m+2)=A(n+1,m+1)+A(n,m+2)$. The inner hypothesis gives $A(n+1,m+1)=\binom{n+m+2}{m+1}-\binom{n+m+2}{m}$ and the outer one gives $A(n,m+2)=\binom{n+m+2}{m+2}-\binom{n+m+2}{m+1}$. Adding $\binom{n+m+3}{m+1}=\binom{n+m+2}{m}+\binom{n+m+2}{m+1}$ (Pascal's rule) gives $A(n+1,m+2)+\binom{n+m+3}{m+1}=\binom{n+m+2}{m+1}+\binom{n+m+2}{m+2}=\binom{n+m+3}{m+2}$.

\emph{Case $m+1$ with $n=m+1$} (diagonal step). By Lemma~\ref{lem:A}(c), $A(m+2,m+2)=A(m+2,m+1)$, and the inner hypothesis gives $A(m+2,m+1)+\binom{2m+3}{m}=\binom{2m+3}{m+1}$. By Pascal's rule and the symmetry $\binom{2m+3}{m+2}=\binom{2m+3}{m+1}$,
\[
  \binom{2m+4}{m+1}=\binom{2m+3}{m}+\binom{2m+3}{m+1},\qquad \binom{2m+4}{m+2}=2\binom{2m+3}{m+1},
\]
so $A(m+2,m+2)+\binom{2m+4}{m+1}=2\binom{2m+3}{m+1}=\binom{2m+4}{m+2}$. (Lean: \texttt{Nat.choose\_succ\_succ'}, \texttt{Nat.choose\_symm\_half}, \texttt{omega}.)
\end{proof}

\begin{lemma}[\decl{A000245_eq}]\label{lem:fact}
Let $j,x\ge0$ with $x+\binom{2j+4}{j}=\binom{2j+4}{j+1}$. Then $\mathrm{A000245}(j+2)=x$.
\end{lemma}

\begin{proof}
From $\binom{2j+4}{j+1}(j+1)!\,(j+3)!=(2j+4)!$ and $\binom{2j+4}{j}\,j!\,(j+4)!=(2j+4)!$ we get, computing in $\mathbb Z$,
\[
  x\,(j+4)!\,(j+1)!=(j+4)\,(2j+4)!-(j+1)\,(2j+4)!=3\,(2j+4)!.
\]
Hence the exact quotient $3(2j+4)!/((j+4)!\,(j+1)!)$, which is $\mathrm{A000245}(j+2)$ by definition, equals $x$.
\end{proof}

\subsection{Proof of Theorem~\ref{thm:main}}
Write $k=K+1$ and $S=\bigl|\{y\in\G(K+1,0):K+1\in y\}\bigr|$; by Lemma~\ref{lem:colsum} we must show $S=\mathrm{A000245}(K)$. The sequences in $\G(K+1,0)$ that do not contain $K+1$ are exactly the elements of $\G(K,0)$, so (\decl{card_split})
\[
  S+N(K,0)=N(K+1,0).
\]
If $K=0$, then $N(1,0)=N(0,0)=1$ (evaluated by \texttt{rfl}), so $S=0=\mathrm{A000245}(0)$. Let $K\ge1$. By Corollary~\ref{cor:N}(a), (c) with $\ell=2$, and (a) again,
\[
  N(K+1,0)=N(K+1,2)=N(K+1,3)+N(K,1)=N(K+1,3)+N(K,0),
\]
so $S=N(K+1,3)=A(K+1,K-1)$. If $K=1$ this is $A(2,0)=1=\mathrm{A000245}(1)$ (Lemma~\ref{lem:A}(a); the value $\mathrm{A000245}(1)=1$ is evaluated by \texttt{decide}). If $K=j+2$, Proposition~\ref{prop:ballot} with $n=j+3$, $m=j$ gives $A(j+3,j+1)+\binom{2j+4}{j}=\binom{2j+4}{j+1}$, and Lemma~\ref{lem:fact} with $x=A(j+3,j+1)$ gives $S=\mathrm{A000245}(j+2)$. \qed

\begin{remark}
For $k\ge2$ the proof shows that $\sum_nT(n,k)=N(k,3)$, the number of sequences in $\G(k,3)$, the empty one included.
\end{remark}

\section{Machine verification}\label{sec:lean}

\subsection*{Files and declarations}
The proof is in \path{lean/Math0/Frontier/OEIS/A388724.lean} in the math0 repository, \url{https://github.com/Foundation-0/math0}. It imports the frozen statement file \path{lean/Math0/Statements/OEIS/A388724.lean} and the audit commands in \path{lean/Math0/Audit.lean}. The main declaration is \decl{Math0.Frontier.A388724.result} \texttt{: statement}, followed by the command
\begin{lstlisting}
#math0_settles Math0.Frontier.A388724.result Math0.Statements.A388724.statement
\end{lstlisting}
All auxiliary declarations in the proof file are private; their names are given with the corresponding results in Section~\ref{sec:proof}. Neither file uses \texttt{sorry}, \texttt{native\_decide}, custom axioms, \texttt{implemented\_by}, \texttt{extern} or \texttt{set\_option}. The proof evaluates only a few small closed terms by computation, such as $N(1,0)=N(0,0)=1$ (\texttt{rfl}) and $\mathrm{A000245}(1)=1$ (\texttt{decide}); it does not use the data checks of the statement file.

\subsection*{Audit}
The command \texttt{\#math0\_settles} (in \path{lean/Math0/Audit.lean}, for statements that math0 formalizes itself) fails unless \decl{statement} is a \texttt{Prop}-valued definition without universe parameters; \decl{statement} and \decl{result} depend on no axioms besides \decl{propext}, \decl{Classical.choice} and \decl{Quot.sound} (this rejects \texttt{sorry}, \texttt{native\_decide} and custom axioms); and the type of \decl{result} is syntactically \decl{statement} or $\neg$\,\decl{statement}, up to metadata, with no unfolding. Running \texttt{lake env lean Math0/Frontier/OEIS/A388724.lean} in \path{lean/} reports
\begin{lstlisting}[basicstyle=\ttfamily\scriptsize]
math0 audit: 'Math0.Statements.A388724.statement' depends on axioms [Classical.choice, Quot.sound, propext]
math0 audit: 'Math0.Frontier.A388724.result' depends on axioms [Classical.choice, Quot.sound, propext]
math0 statement check: 'Math0.Frontier.A388724.result' proves 'Math0.Statements.A388724.statement'
\end{lstlisting}

\subsection*{Toolchain and reproduction}
We use Lean~4 \cite{Lean4} \texttt{v4.33.1} and Mathlib \cite{Mathlib} \texttt{v4.33.1}, pinned in \path{lean/}; the repository also pins formal-conjectures \cite{FC} at commit \texttt{c77fdd4}, which this result does not use. To reproduce, run \path{scripts/verify.sh} from the repository root: it fetches the Mathlib cache, builds all math0 results with \texttt{lake build} (re-running every audit command) and checks that the pinned \texttt{formal\_conjectures} and \texttt{mathlib} checkouts are unmodified.

{\footnotesize
\begin{thebibliography}{99}

\bibitem{ABEK}
M.~Archibald, A.~Blecher, S.~Elizalde, and A.~Knopfmacher, \emph{Subdiagonal and superdiagonal partitions}, Afrika Matematika \textbf{36} (2025), no.~2, Article~77, \url{https://doi.org/10.1007/s13370-025-01282-0}.

\bibitem{Deutsch}
E.~Deutsch, \emph{Dyck path enumeration}, Discrete Math. \textbf{204} (1999), 167--202, \url{https://doi.org/10.1016/S0012-365X(98)00371-9}.

\bibitem{FC}
M.~Firsching, P.~Lezeau, S.~Mercuri, M.~Z. Horv\'ath, Y.~Dillies, C.~S\"onne, E.~Wieser, F.~Zhang, T.~Hubert, B.~Ag\"uera y Arcas, and P.~Kohli, \emph{Formal Conjectures: An Open and Evolving Benchmark for Verified Discovery in Mathematics}, arXiv:2605.13171, 2026. Repository: \url{https://github.com/google-deepmind/formal-conjectures}.

\bibitem{Mathlib}
The mathlib Community, \emph{The Lean mathematical library}, in: Proceedings of the 9th ACM SIGPLAN International Conference on Certified Programs and Proofs (CPP 2020), ACM, 2020, 367--381.

\bibitem{Lean4}
L.~de Moura and S.~Ullrich, \emph{The Lean 4 theorem prover and programming language}, in: Automated Deduction -- CADE 28, Lecture Notes in Comput. Sci. \textbf{12699}, Springer, 2021, 625--635.

\bibitem{A000245}
The OEIS Foundation Inc., \emph{$a(n)=3(2n)!/((n+2)!\,(n-1)!)$}, entry A000245 (N.~J.~A. Sloane) in The On-Line Encyclopedia of Integer Sequences, \url{https://oeis.org/A000245}.

\bibitem{A009766}
The OEIS Foundation Inc., \emph{Catalan's triangle $T(n,k)$ (read by rows)}, entry A009766 (W.~Meeussen) in The On-Line Encyclopedia of Integer Sequences, \url{https://oeis.org/A009766}.

\bibitem{A033184}
The OEIS Foundation Inc., \emph{Catalan triangle A009766 transposed}, entry A033184 (C.~G. Bower) in The On-Line Encyclopedia of Integer Sequences, \url{https://oeis.org/A033184}.

\bibitem{A388724}
The OEIS Foundation Inc., \emph{Triangle read by rows where $T(n,k)$ is the number of strictly superdiagonal reversed integer partitions of $n$ with maximum $k$}, entry A388724 (G.~Wiseman, 2025) in The On-Line Encyclopedia of Integer Sequences, \url{https://oeis.org/A388724}; revision \#12, October~6, 2025.

\end{thebibliography}
}

\end{document}
